\documentclass[11pt]{article}
\usepackage[margin=1in]{geometry}
\usepackage{amsmath,amssymb,amsthm,listings}
\usepackage[colorlinks=true,urlcolor=blue,linkcolor=blue]{hyperref}
\lstset{basicstyle=\ttfamily\small,columns=fullflexible,
  keepspaces=true,showstringspaces=false,breaklines=true}
\newtheorem{theorem}{Theorem}
\title{An exact counterexample to cosine-sum positivity\\
\large MathDB 369172}
\author{Principia Math In Progress}
\date{11 September 2026}
\begin{document}
\maketitle

\section*{Task 1. Resolve the positivity conjecture}
For an integer $d\geq2$, choose the square-root branch equal to $1$ at
the origin and define
\[
 \sqrt{1+z+\cdots+z^d}=\sum_{k\geq0}b_k^{(d)}z^k,\qquad
 \mathcal T_n^{(d)}(t)=\sum_{k=0}^n b_k^{(d)}\cos(kt).
\]
Kovalev's Conjecture 1.3~\cite{K} asserts that these cosine sums are
positive for every $d\geq2$, $n\geq0$ and real $t$. It is false.

\begin{theorem}
Let $d=43$, $n=1800$ and $t=\arccos(99/100)$. Then
\[
 -\frac{61}{10000}
 <\mathcal T_{1800}^{(43)}(t)
 <-\frac{6099}{1000000}<0.
\]
\end{theorem}

\begin{proof}
All parts of the counterexample can be checked with integers.
Put $q=d+1=44$ and
\[
 f(z)=\sqrt{\frac{1-z^q}{1-z}}=\sum_{k\geq0}b_kz^k.
\]
Logarithmic differentiation gives
\[
 2(1-z)(1-z^q)f'(z)
   =\bigl(1-qz^{q-1}+(q-1)z^q\bigr)f(z).
\]
Comparing coefficients, with negative-index coefficients set to zero,
yields
\[
 2k b_k=(2k-1)b_{k-1}
       +(2k-3q)b_{k-q}+(3q+1-2k)b_{k-q-1}.
\]
Consequently the scaled coefficients $B_k=4^k b_k$ satisfy $B_0=1$ and
\begin{align}
 2kB_k={}&4(2k-1)B_{k-1}
   +4^{44}(2k-132)B_{k-44}\notag\\
   &+4^{45}(133-2k)B_{k-45}.
 \label{eq:B}
\end{align}
They are integers. For example, this follows independently from the
binomial expansions:
\begin{equation}
 B_k=\binom{2k}{k}
 -\sum_{j=1}^{\lfloor k/44\rfloor}
 \frac{\binom{2j}{j}}{2j-1}\,
 4^{43j}\binom{2(k-44j)}{k-44j}.
 \label{eq:independent}
\end{equation}
Indeed $\binom{2j}{j}/(2j-1)=2\operatorname{Cat}_{j-1}$ is an integer.
This identity also fixes the branch and normalization unambiguously.

Define integers $C_0=1$, $C_1=99$ and
\[
 C_k=198C_{k-1}-10000C_{k-2}\quad(k\geq2).
\]
The cosine addition formula gives $C_k=100^k\cos(kt)$.
Finally put
\[
 A_0=1,\qquad A_k=400A_{k-1}+B_kC_k.
\]
Induction shows that
\[
 \frac{A_k}{400^k}
   =\sum_{j=0}^k\frac{B_jC_j}{400^j}
   =\mathcal T_k^{(43)}(t).
\]
The following complete integer calculation proves
\[
 10000A_{1800}+61\cdot400^{1800}>0,\qquad
 -6099\cdot400^{1800}-1000000A_{1800}>0.
\]
It uses only standard Python integer arithmetic. Every division in
\eqref{eq:B} is checked for zero remainder.

\begin{lstlisting}[language=Python]
B = [1]
def at(k):
    return B[k] if k >= 0 else 0
for k in range(1, 1801):
    u = 4*(2*k-1)*at(k-1)
    u += 4**44*(2*k-132)*at(k-44)
    u += 4**45*(133-2*k)*at(k-45)
    v, r = divmod(u, 2*k)
    if r != 0:
        raise ArithmeticError(k)
    B.append(v)
C = [1, 99]
for k in range(2, 1801):
    C.append(198*C[-1] - 10000*C[-2])
A = 1
for k in range(1, 1801):
    A = 400*A + B[k]*C[k]
D = 400**1800
if not (10000*A + 61*D > 0
        and -6099*D - 1000000*A > 0):
    raise ArithmeticError("certificate failed")
print("Both strict bounds verified.")
\end{lstlisting}
Dividing by the positive denominator proves the theorem.
\end{proof}

\paragraph{Independent checks.}
The accompanying \texttt{verify\_counterexample.py} additionally checks
all 1801 coefficients against~\eqref{eq:independent}, all cosine values
against the real parts of $(99+i\sqrt{199})^k$, and the final sum in the
reverse order. These checks use integers throughout. The scaled numerator
has 15552 bits in absolute value; the denominator has 15559 bits.
The saved certificate records both strict bounds and the exact numerator
in hexadecimal.

\paragraph{Scope.}
This refutes Conjecture 1.3 and MathDB 369172. A negative real part at
one boundary point does not exhibit a zero of the Taylor polynomial.
The distinct zero-free Conjecture 1.2, recorded as MathDB 369171,
therefore remains unresolved here. No minimality of $d$ or $n$ is claimed.

\section*{Task 2. Determine the parameter range for positivity}
Classify the integers $d\geq2$ for which
$\mathcal T_n^{(d)}(t)>0$ holds for every $n$ and real $t$.
In particular, settle the source's $d=2$ positivity question.

\section*{Task 3. Determine the first failures}
For each $d$ admitting a negative cosine sum, determine the least
truncation degree at which positivity fails and certify the relevant
angles or intervals of angles.

\section*{Task 4. Prove a boundary scaling law}
At each nontrivial $(d+1)$st root of unity, derive a uniform local
asymptotic formula for the Taylor sections at angular distance of
order $1/n$. Use it to identify parameter ranges with infinitely many
negative cosine sums.

\section*{Task 5. Resolve the separate zero-free conjecture}
Determine whether every Taylor section of
$\sqrt{1+z+\cdots+z^d}$, $d>2$, is zero-free in the closed unit disk.
An affirmative proof must accommodate the negative real parts proved here.

\begin{thebibliography}{9}
\bibitem{K}
L. V. Kovalev,
\emph{Sharp bounds for some segments of bounded power series},
arXiv:2507.04544v1 (2025), Conjectures 1.2 and 1.3, p.~2.
\url{https://arxiv.org/abs/2507.04544}.
\end{thebibliography}
\end{document}
